% !TITLE 诗经·周南·桃夭
% !AUTHOR 无名氏
% !DYNASTY 先秦
% !GENRE 四言诗
% !ORDER 6

\begin{poem}
桃之夭夭\textsuperscript{1}，灼灼\textsuperscript{2}其华。
之子于归\textsuperscript{3}，宜\textsuperscript{4}其室家。

桃之夭夭，有蕡\textsuperscript{5}其实。
之子于归，宜其家室。

桃之夭夭，其叶蓁蓁\textsuperscript{6}。
之子于归，宜其家人。
\end{poem}

\begin{annotations}
\notetext{1}{夭夭：树木柔嫩茂盛的样子。这里也用来比喻新娘的年轻美丽。}
\notetext{2}{灼灼：花开得红火鲜艳的样子。华，同“花”。}
\notetext{3}{之子于归：之子，这个女子。于归，古代对女子出嫁的雅称，意思是回到最终的归宿（夫家）。}
\notetext{4}{宜：使……和顺适宜，这里指相处和睦。}
\notetext{5}{蕡：蕡（fén），果实繁多、又大又圆的样子。}
\notetext{6}{蓁蓁：蓁蓁（zhēn zhēn），树叶繁密茂盛的样子。}
\end{annotations}

\begin{appreciation}
这是一首祝贺女子出嫁的诗，也是中国文学史上最早、最美的婚礼贺词之一。

闻一多说，第一个想出用花来比喻美女的是天才，第二个想出来的是庸才，而第三个想出来的是蠢才。这首诗的作者就是他口中的天才。从这以后，“桃花”就成了美女的代称。一个人女人缘好，就说他很有桃花运。

这首诗的每一章都以桃树的不同部分起兴，祝福这位年轻美丽的女子出嫁后家庭美满、婚姻和谐。每一章之间的内容变化不多，体现了《诗经》重章叠句的特点，就像歌词一样朗朗上口。的确，这些来自国风的诗，正是从民歌中采集而来。借助这首诗，我们仿佛又能看到一场祖先的婚礼。岁月流逝，但真挚的祝福却始终如一。

也祝你桃之夭夭，灼灼其华。
\end{appreciation}
